% Supplementary material, submitted as a separate PDF.
%
% TDSC counts pages beyond 12 as overlength and caps the paper at 18; supplementary files have no
% page limit. The full hybrid proof lives here so the main text can carry the proof sketch and the
% measurements.
%
% Build main.tex FIRST: the xr package reads main.aux so that references to Theorem 1, the
% definitions and the evaluation sections resolve to the numbers the main paper gives them.
%
%   latexmk -pdf main.tex && latexmk -pdf supplement.tex

\documentclass[10pt,journal,compsoc]{IEEEtran}

\usepackage[T1]{fontenc}
\usepackage[utf8]{inputenc}
\usepackage{amsmath,amssymb}
\let\proof\relax
\let\endproof\relax
\usepackage{amsthm}
\usepackage{booktabs}
\usepackage{graphicx}
\usepackage{microtype}
\usepackage[hidelinks]{hyperref}
\usepackage[capitalise,nameinlink]{cleveref}

% Resolve \ref and \cref against the main paper's numbering.
\usepackage{xr}
\externaldocument{main}

\newcommand{\advhist}{\ensuremath{\mathrm{Adv}_{\mathrm{hist}}}}
\newcommand{\advpost}{\ensuremath{\mathrm{Adv}_{\mathrm{post}}}}
\newcommand{\AdvH}{\ensuremath{\mathsf{Adv}^{\mathrm{H}}}}
\newcommand{\AdvR}{\ensuremath{\mathsf{Adv}^{\mathrm{R}}}}
\newcommand{\negl}{\ensuremath{\mathsf{negl}}}
% Every number the prose quotes, generated from the results (python -m eval.exp2.numbers).
% Generated by `python -m eval.exp2.numbers` from python/eval/results/*.json.
% Do not edit: re-run the generator. Every macro is a number the prose quotes.
\newcommand{\RDeleteN}{100}
\newcommand{\RDeleteNamedP}{0}
\newcommand{\RDeleteNamedRS}{99}
\newcommand{\REeighteenAUC}{0.80~[0.70, 0.90]}
\newcommand{\REeighteenAUCPt}{0.80}
\newcommand{\REeighteenMemberWrite}{0.30}
\newcommand{\REeighteenModel}{qwen2.5:7b-instruct}
\newcommand{\REeighteenNonmemberWrite}{0.90}
\newcommand{\REelevenFiqa}{0.384}
\newcommand{\REelevenFiqaSize}{57{,}638}
\newcommand{\REelevenGroupedExhFiqa}{0.384}
\newcommand{\REelevenGroupedExhScifact}{0.738}
\newcommand{\REelevenGroupedFiqa}{0.224}
\newcommand{\REelevenGroupedSFiqa}{0.329}
\newcommand{\REelevenGroupedSScifact}{0.658}
\newcommand{\REelevenGroupedScifact}{0.435}
\newcommand{\REelevenIvfGainMax}{+0.045}
\newcommand{\REelevenIvfGainMin}{+0.023}
\newcommand{\REelevenProbeCostMax}{0.080}
\newcommand{\REelevenRefusals}{0}
\newcommand{\REelevenWorstQuant}{0.001}
\newcommand{\REfifteenARIGrouped}{0.59}
\newcommand{\REfifteenARISfour}{0.13}
\newcommand{\REfifteenARISone}{0.94}
\newcommand{\REfifteenARIStwo}{0.66}
\newcommand{\REfifteenExactGrouped}{0.85}
\newcommand{\REfifteenExactSfour}{0.00}
\newcommand{\REfifteenExactSone}{0.99}
\newcommand{\REfifteenExactStwo}{0.00}
\newcommand{\REfifteenQueries}{16}
\newcommand{\REfifteenRepeats}{20}
\newcommand{\REfifteenSearches}{10}
\newcommand{\REfifteenThreshGrouped}{0.85~[0.82, 0.88]}
\newcommand{\REfifteenThreshSfour}{0.59~[0.57, 0.60]}
\newcommand{\REfifteenThreshSone}{0.98~[0.97, 1.00]}
\newcommand{\REfifteenThreshStwo}{0.98~[0.97, 0.99]}
\newcommand{\REfiveInstances}{1{,}000}
\newcommand{\REfiveRegionFirstFree}{0.51~[0.48, 0.55]}
\newcommand{\REfiveShipped}{0.50~[0.47, 0.54]}
\newcommand{\REfiveSlotFirstFree}{1.00~[1.00, 1.00]}
\newcommand{\REfiveSlotRandom}{0.49~[0.45, 0.53]}
\newcommand{\REfourAware}{+0.42}
\newcommand{\REfourAwareMajority}{0.004~[-0.015, 0.023]}
\newcommand{\REfourMajority}{0.50}
\newcommand{\REfourNaive}{\ensuremath{-}0.02}
\newcommand{\REfourShuffled}{+0.37}
\newcommand{\REninePaddedAUC}{0.500~[0.500, 0.500]}
\newcommand{\REninePaddedAUCPt}{0.50}
\newcommand{\REninePaddedFull}{1.000}
\newcommand{\REninePaddedOpen}{1.000}
\newcommand{\REnineSilentAUC}{1.000~[1.000, 1.000]}
\newcommand{\REnineSilentAUCPt}{1.00}
\newcommand{\REnineSilentFull}{0.070}
\newcommand{\REnineSilentOpen}{1.000}
\newcommand{\REnineUtilM}{100.0}
\newcommand{\REnineUtilP}{23.8}
\newcommand{\REnineteenAddedMOne}{959}
\newcommand{\REnineteenAddedPOne}{1099}
\newcommand{\REnineteenAmpMOne}{110}
\newcommand{\REnineteenAmpPOne}{13}
\newcommand{\REnineteenMBlobs}{1.0}
\newcommand{\REnineteenPrecPFour}{0.058~[0.057, 0.058]}
\newcommand{\REnineteenPrecPOne}{0.207~[0.205, 0.209]}
\newcommand{\REnineteenPrecRatio}{13}
\newcommand{\REnineteenRewrittenMOne}{9}
\newcommand{\REnineteenRewrittenPOne}{85}
\newcommand{\REoneGrowingMax}{0.097~[0.092, 0.103]}
\newcommand{\REoneMMax}{0.001~[0.000, 0.002]}
\newcommand{\REonePMax}{0.097~[0.085, 0.108]}
\newcommand{\REoneRSMax}{0.065~[0.060, 0.072]}
\newcommand{\REsevenAllExact}{18}
\newcommand{\REsevenExactCIGrouped}{0.025~[0.016, 0.034]}
\newcommand{\REsevenExactCISone}{0.963~[0.897, 1.000]}
\newcommand{\REsevenExactCIStwo}{0.564~[0.469, 0.664]}
\newcommand{\REsevenExactGrouped}{0.025}
\newcommand{\REsevenExactSone}{0.963}
\newcommand{\REsevenExactStwo}{0.564}
\newcommand{\REsevenFaqSone}{0.887}
\newcommand{\REsevenFaqStwo}{0.393}
\newcommand{\REsevenNeighGrouped}{0.172}
\newcommand{\REsevenNeighSone}{0.965}
\newcommand{\REsevenNeighStwo}{1.000}
\newcommand{\REsevenRegionsGrouped}{2}
\newcommand{\REsevenRegionsSone}{62}
\newcommand{\REsevenRegionsStwo}{36}
\newcommand{\REsevenRepeats}{20}
\newcommand{\REseventeenCandFour}{17.3}
\newcommand{\REseventeenCandOne}{4.8}
\newcommand{\REseventeenCandSixteen}{46.3}
\newcommand{\REseventeenCandTwo}{9.3}
\newcommand{\REseventeenComposedFour}{\ensuremath{-}0.079}
\newcommand{\REseventeenComposedOne}{\ensuremath{-}0.035}
\newcommand{\REseventeenComposedSixteen}{\ensuremath{-}0.096}
\newcommand{\REseventeenComposedTwo}{\ensuremath{-}0.057}
\newcommand{\REseventeenOracle}{\ensuremath{-}0.003}
\newcommand{\REsixChroma}{\ensuremath{-}0.011~[\ensuremath{-}0.025, +0.004]}
\newcommand{\REsixSeeds}{10}
\newcommand{\REsixWorstHi}{+0.004}
\newcommand{\REsixteenPgAcc}{0.243}
\newcommand{\REsixteenPgAdv}{0.118~[0.094, 0.143]}
\newcommand{\REsixteenPgAdvPt}{+0.12}
\newcommand{\REsixteenPgChance}{0.125}
\newcommand{\REsixteenPgSeeds}{20}
\newcommand{\REtenCrossHi}{30{,}000}
\newcommand{\REtenCrossLo}{10{,}000}
\newcommand{\REtenIvfSearchLarge}{1.78}
\newcommand{\REtenIvfSpeedupMax}{18.3}
\newcommand{\REtenIvfSpeedupMin}{3.8}
\newcommand{\REtenMSearchLarge}{23.19}
\newcommand{\REtenNLarge}{100{,}000}
\newcommand{\REtenNSmall}{1{,}000}
\newcommand{\REtenPSearchLarge}{6.67}
\newcommand{\REtenPutMLarge}{32.2}
\newcommand{\REtenPutPLarge}{46.3}
\newcommand{\REtenResidentLarge}{544}
\newcommand{\REtenSpeedupLarge}{0.29}
\newcommand{\REtenSpeedupSmall}{9.1}
\newcommand{\REtenUnlockIvfLarge}{6.7}
\newcommand{\REtenUnlockMLarge}{0.8}
\newcommand{\REthirtyAUC}{0.50~[0.50, 0.50]}
\newcommand{\REthirtyMemberWrite}{1.00}
\newcommand{\REthirtyNonmemberWrite}{1.00}
\newcommand{\REthirtyoneGemma}{0.62~[0.56, 0.69]}
\newcommand{\REthirtyoneGemmaMem}{0.72}
\newcommand{\REthirtyoneGemmaNon}{0.97}
\newcommand{\REthirtyoneLlama}{0.48~[0.44, 0.51]}
\newcommand{\REthirtyoneLlamaMem}{0.98}
\newcommand{\REthirtyoneLlamaNon}{0.93}
\newcommand{\REthirtyoneMistral}{0.78~[0.72, 0.84]}
\newcommand{\REthirtyoneMistralMem}{0.43}
\newcommand{\REthirtyoneMistralNon}{1.00}
\newcommand{\REthirtyoneQwen}{0.74~[0.67, 0.81]}
\newcommand{\REthirtyoneQwenKOne}{0.742}
\newcommand{\REthirtyoneQwenKThree}{0.745}
\newcommand{\REthirtyoneQwenMem}{0.47}
\newcommand{\REthirtyoneQwenNon}{0.95}
\newcommand{\REthirtythreeAgentExposure}{1.00}
\newcommand{\REthirtythreeAgentFlushes}{3.0}
\newcommand{\REthirtythreeAgentRatio}{4.7}
\newcommand{\REthirtythreeFastExposure}{1.00}
\newcommand{\REthirtythreeFastFlushes}{13.9}
\newcommand{\REthirtythreeFastRatio}{1.4}
\newcommand{\REthirtytwoFour}{0.067~[0.043, 0.092]}
\newcommand{\REthirtytwoOne}{0.187~[0.155, 0.219]}
\newcommand{\REthirtytwoSeeds}{10}
\newcommand{\REthirtytwoSixteen}{0.000~[0.000, 0.000]}
\newcommand{\REthreeLlmWriteMember}{0.74}
\newcommand{\REthreeLlmWriteNonmember}{0.97}
\newcommand{\REthreeMLlmAUC}{0.50}
\newcommand{\REthreeModel}{qwen2.5:7b-instruct}
\newcommand{\REthreePLlmAUC}{0.62~[0.57, 0.67]}
\newcommand{\REtwentyMContent}{0.976~[0.951, 0.995]}
\newcommand{\REtwentyMOnDemand}{1.000~[1.000, 1.000]}
\newcommand{\REtwentyMRegions}{0.490~[0.456, 0.523]}
\newcommand{\REtwentyMSweep}{0.505~[0.475, 0.538]}
\newcommand{\REtwentyMSweepPt}{0.50}
\newcommand{\REtwentyPShipped}{1.000~[1.000, 1.000]}
\newcommand{\REtwentyPShippedPt}{1.00}
\newcommand{\REtwentyeightGrowOneMin}{3}
\newcommand{\REtwentyeightGrowOneMsgs}{40}
\newcommand{\REtwentyeightPOneMin}{26}
\newcommand{\REtwentyeightPOneMsgs}{320}
\newcommand{\REtwentyeightPTwoMin}{3}
\newcommand{\REtwentyeightPTwoMsgs}{40}
\newcommand{\REtwentyeightSecHi}{4.9}
\newcommand{\REtwentyeightSecLo}{3.5}
\newcommand{\REtwentyfiveDiskPut}{1.6}
\newcommand{\REtwentyfiveDiskRAM}{0.6}
\newcommand{\REtwentyfiveDiskSearch}{13.5}
\newcommand{\REtwentyfiveDiskUnlock}{8.2}
\newcommand{\REtwentyfiveMemPut}{1.3}
\newcommand{\REtwentyfiveMemRAM}{544.3}
\newcommand{\REtwentyfiveMemSearch}{0.9}
\newcommand{\REtwentyfiveMemUnlock}{6.6}
\newcommand{\REtwentyfivePPut}{42.4}
\newcommand{\REtwentyfivePRAM}{0.0}
\newcommand{\REtwentyfivePSearch}{6.5}
\newcommand{\REtwentyfivePUnlock}{0.7}
\newcommand{\REtwentyfourIvfOwn}{0.657}
\newcommand{\REtwentyfourIvfPub}{0.672}
\newcommand{\REtwentyfourMOwn}{0.703}
\newcommand{\REtwentyfourMPub}{0.708}
\newcommand{\REtwentyfourPOwn}{0.665}
\newcommand{\REtwentyfourPPub}{0.441}
\newcommand{\REtwentyfourRefusedOwn}{0}
\newcommand{\REtwentyfourRefusedPub}{2{,}297}
\newcommand{\REtwentyfourShare}{9.9}
\newcommand{\REtwentynineMAll}{0.63~[0.52, 0.73]}
\newcommand{\REtwentynineMAllMem}{540}
\newcommand{\REtwentynineMAllNon}{539}
\newcommand{\REtwentynineMAllPt}{0.63}
\newcommand{\REtwentynineMStore}{0.47~[0.37, 0.56]}
\newcommand{\REtwentynineMStoreMem}{652}
\newcommand{\REtwentynineMStoreNon}{652}
\newcommand{\REtwentynineMStorePt}{0.47}
\newcommand{\REtwentyninePStore}{0.69~[0.60, 0.78]}
\newcommand{\REtwentyninePStoreMem}{831}
\newcommand{\REtwentyninePStoreNon}{1528}
\newcommand{\REtwentyninePStorePt}{0.69}
\newcommand{\REtwentyninePerClass}{60}
\newcommand{\REtwentynineQdrant}{0.83~[0.76, 0.90]}
\newcommand{\REtwentynineQdrantMem}{2}
\newcommand{\REtwentynineQdrantNon}{7}
\newcommand{\REtwentynineQdrantPt}{0.83}
\newcommand{\REtwentynineWindow}{16}
\newcommand{\REtwentyoneDriveYears}{2{,}718}
\newcommand{\REtwentyonePLockMiB}{78}
\newcommand{\REtwentyoneRatio}{1{,}562}
\newcommand{\REtwentyoneRegionGBDay}{945}
\newcommand{\REtwentyoneRegionMiB}{5.2}
\newcommand{\REtwentyoneResidualHi}{35}
\newcommand{\REtwentyoneResidualLo}{25}
\newcommand{\REtwentyoneResidualSustain}{0}
\newcommand{\REtwentyoneSessions}{40}
\newcommand{\REtwentyoneSweepGBDay}{0.60}
\newcommand{\REtwentyoneSweepKiB}{68.4}
\newcommand{\REtwentyoneSweepMs}{8.1}
\newcommand{\REtwentysevenFull}{+0.405}
\newcommand{\REtwentysevenQuarter}{+0.216}
\newcommand{\REtwentysevenTenth}{+0.103}
\newcommand{\REtwentysevenZero}{\ensuremath{-}0.019}
\newcommand{\REtwentysixNfMIvf}{0.318}
\newcommand{\REtwentysixNfP}{0.291}
\newcommand{\REtwentysixNfQuant}{\ensuremath{-}0.000}
\newcommand{\REtwentysixSciMIvf}{0.652}
\newcommand{\REtwentysixSciP}{0.604}
\newcommand{\REtwentysixSciQuant}{\ensuremath{-}0.000}
\newcommand{\REtwentythreeMWorst}{100}
\newcommand{\REtwentythreePBest}{24}
\newcommand{\REtwentythreePSkewHigh}{6}
\newcommand{\REtwentythreePSkewLow}{9}
\newcommand{\REtwentythreePWorst}{6}
\newcommand{\REtwentytwoGroupAmpEight}{2.0}
\newcommand{\REtwentytwoGroupAmpFour}{1.0}
\newcommand{\REtwentytwoGroupAmpSixteen}{4.0}
\newcommand{\REtwentytwoGroupAmpThirtyTwo}{8.0}
\newcommand{\REtwentytwoGroupAmpTwo}{0.5}
\newcommand{\REtwentytwoGroupEight}{0.834~[0.800, 0.866]}
\newcommand{\REtwentytwoGroupFour}{0.920~[0.896, 0.941]}
\newcommand{\REtwentytwoGroupSixteen}{0.691~[0.637, 0.743]}
\newcommand{\REtwentytwoGroupThirtyTwo}{0.459~[0.393, 0.509]}
\newcommand{\REtwentytwoGroupTwo}{0.959~[0.943, 0.973]}
\newcommand{\REtwentytwoSpreadEight}{0.186~[0.173, 0.197]}
\newcommand{\REtwentytwoSpreadFour}{0.574~[0.558, 0.590]}
\newcommand{\REtwentytwoSpreadOne}{0.966~[0.932, 0.991]}
\newcommand{\REtwentytwoSpreadSix}{0.344~[0.330, 0.359]}
\newcommand{\REtwentytwoSpreadThree}{0.820~[0.803, 0.836]}
\newcommand{\REtwentytwoSpreadTwo}{0.978~[0.970, 0.985]}
\newcommand{\REtwoBurstMax}{0.065~[0.054, 0.074]}
\newcommand{\REtwoBurstRatio}{9.0}
\newcommand{\REtwoGrowingMax}{0.188~[0.173, 0.203]}
\newcommand{\REtwoSingleMaxHi}{0.018}
\newcommand{\RLedgerPermBits}{4.9}
\newcommand{\RLedgerPermBitsMax}{6}
\newcommand{\RSearchPositionsRS}{292}
\newcommand{\RSearchReadsP}{0}
\newcommand{\RSearchReadsRS}{2{,}000}

% Run-in paragraph heads unnumbered, as in the main paper.
\setcounter{secnumdepth}{1}

\crefname{section}{Section}{Sections}

\begin{document}

\title{Supplementary Material for\\
``Historical Confidentiality for Encrypted Agent Memory:\\ Dual-Epoch Security and the Channels
That Defeat It''}

\author{Shreyas~Fegade%
\IEEEcompsocitemizethanks{\IEEEcompsocthanksitem S. Fegade is an independent researcher.\protect\\
E-mail: shreyasf@icloud.com}}

\markboth{Supplementary Material}{Fegade: Historical Confidentiality for Encrypted Agent Memory}

\maketitle

\noindent
This document contains the full proof of Theorem~\ref{thm:main}, whose statement, conditions and
proof sketch appear in the main paper. Numbering of theorems, definitions, propositions and
sections follows the main paper throughout.

\section{Proof of Theorem~\ref{thm:main}}
\label{supp:proof}

Let $\Pi$ satisfy S1 to S4, let $F$ be the pseudorandom function used to derive record tags and
the per-flush random choices, and let $\mathsf{AE}$ be the authenticated encryption scheme. Let
$\mathcal{A}$ be a probabilistic polynomial-time adversary in $\mathbf{G}^{\mathrm{H}}$ that
submits two equal-length histories --- admissible or not --- and issues polynomially many post-$\tau$
operations. The proof proceeds through three games and then exhibits the simulator of
\cref{def:dualepoch}.

\paragraph{Cryptographic model} $\mathsf{AE}$ is
XChaCha20-Poly1305, used as an authenticated cipher that is IND\$-CPA (ciphertexts are
indistinguishable from uniform strings of the same length) and INT-CTXT. Nonces are $192$-bit and
drawn uniformly at random per encryption; the width makes a collision negligible ($q^2/2^{193}$ over
$q$ encryptions) and removes any nonce counter from disk. A header's associated data is a
label naming the engine and format version together with the header's position (region and slot in
Mnemos-P, absolute slot in Mnemos-M), so a header moved to another slot fails authentication; a
payload's associated data is its record's identifier tag, which binds it to its header. Headers,
payloads, the manifest, the table and the journal are sealed under distinct derived keys. The manifest and the partition
table are themselves $\mathsf{AE}$ ciphertexts of fixed length under their own derived keys; the
partition table is a ciphertext, not a PRF image, and is included in the snapshot as one of the
fixed-length uniform strings of game $\mathrm{G}_2$. All working keys are derived from the master by
HKDF, treated as a PRF with independent labels per key, so the derived keys are independent. The
simulator is given the public structural parameters $(B, C, d, \text{block size}, w, \Delta, k)$ and
the leakage $\mathcal{L}$; it is \emph{not} given $b$, the histories, or any plaintext. Adaptivity is
permitted on physical locations the adversary has already observed, which the next paragraph treats.

\paragraph{On adaptivity} $\mathcal{A}$ chooses each post-$\tau$ operation after seeing the leakage
of every earlier one. This does not disturb the argument. Each hybrid step below replaces a
primitive with an idealized version, uniformly over the whole game, rather than fixing the
operation sequence in advance; the reduction answers $\mathcal{A}$'s queries as they arrive and its
advantage is unchanged. The final step shows the view is a function of the leakage and public
randomness alone, which holds pointwise for every sequence $\mathcal{A}$ might produce and
therefore holds for an adaptively chosen one.

\paragraph{Game $\mathrm{G}_0$} The real game $\mathbf{G}^{\mathrm{H}}_{\Pi,\mathcal{A}}$.

\paragraph{Game $\mathrm{G}_1$} As $\mathrm{G}_0$, except that every evaluation of $F$ under a
store key is replaced by a lazily sampled random function. This affects the record tags, the
keyed logical-to-physical partition permutation, the free-slot choice within a rewritten region,
the choice of uniformly random padding regions in a flush, and the per-call probe randomization.
A distinguisher between $\mathrm{G}_0$ and $\mathrm{G}_1$ yields a distinguisher against $F$, so
$|\Pr[\mathrm{G}_0 \Rightarrow 1] - \Pr[\mathrm{G}_1 \Rightarrow 1]| \le
\mathsf{Adv}^{\mathrm{prf}}_{F}(\lambda)$ up to a factor for the number of derived keys.

\paragraph{Game $\mathrm{G}_2$} As $\mathrm{G}_1$, except that every ciphertext written to disk is
replaced by a uniformly random string of the same length. This includes each slot's header and
payload ciphertexts, the manifest, and the partition table, and every re-encryption performed when
a region is rewritten during a flush. The challenger keeps the plaintexts of live records in its
own state, so no stored ciphertext is decrypted. Each ciphertext is produced under a fresh
uniformly random nonce and the number of encryptions is polynomial, so a reduction answering every
encryption with an IND\$-CPA challenge oracle gives
$|\Pr[\mathrm{G}_1 \Rightarrow 1] - \Pr[\mathrm{G}_2 \Rightarrow 1]| \le
\mathsf{Adv}^{\mathrm{ind\$}}_{\mathsf{AE}}(\lambda)$.

\paragraph{Independence of the view in $\mathrm{G}_2$ from $b$} The view consists of the snapshot
at $\tau$ and the leakage of each post-$\tau$ operation. Each component is considered in turn.

\begin{enumerate}
\item \emph{File lengths.} By S1 every file has a length determined by the public parameters, so
the snapshot's shape is independent of $b$.
\item \emph{Snapshot bytes.} By S2 and the change made in $\mathrm{G}_2$, every slot of the
partition file, whether it holds a live record or is free, is a uniformly random string of the
slot length, and the manifest and table are uniformly random strings of fixed length. The snapshot
is therefore a tuple of independent uniform strings whose count and lengths are public parameters,
identically distributed under $b=0$ and $b=1$.
\item \emph{Flush leakage.} By S3 a flush writes exactly $w+1$ whole regions. The regions written
are: one per staged record, being that record's physical region; uniformly random regions to reach
$w$; and one region fixed by a counter that depends only on the flush index. A staged record's
physical region is the relabeling of the nearest-centroid partition of its embedding, a function of
the embedding, the partition table and the relabeling, all independent of $b$ by S4 and
$\mathrm{G}_1$. The uniform and counter choices are independent of $b$. The written bytes are
uniform by item 2. So a flush reveals only a multiset of region indices with a distribution
independent of $b$.
\item \emph{Deletion and update leakage.} By S3 a deletion writes nothing at the time of the call
and records a tombstone in the manifest, which is rewritten wholesale every flush regardless of
content; an update is a deletion followed by a fresh write. The physical removal of a tombstoned
record occurs only when the counter-selected region of a later flush is rewritten, which is already
covered by item 3 and whose region choice is content-independent. Hence deleting or updating a
cohort-H record adds nothing to the view that depends on which region held it.
\item \emph{Search leakage.} A search reads the regions its probe set names. The probe set is a
function of the query, the partition table, the relabeling and fresh per-call randomness, none of
which depends on $b$. Ranking and payload decryption occur in memory over the region images already
read, so no further disk access is issued.
\item \emph{Refusals.} No admissibility assumption is needed: neither history is required to stay
below capacity, and the case is closed by S3's padded refusal. Consider a post-$\tau$
$\mathsf{put}$ whose target partition is full under one history and not under the other. Under both,
S3 requires the store to rewrite the intended region: an accepted write re-seals the region with the
new record among its slots, and a refused write re-seals the same region unchanged. By item~2 every
slot of a rewritten region is a uniform string of the slot length whether or not it holds a live
record, so the two rewrites are identically distributed, and the region index written is the same
function of the (relabeled) target partition in both cases. The observation extends to a
$\mathsf{put}$ that overflows a partition \emph{within} $\sigma_b$, before $\tau$: the refused writes
rewrite their regions, and by item~2 the snapshot does not reveal a region's occupancy, so a history
that overflows a partition and one that does not produce identically distributed snapshots. Hence
the refusal rule is exercised freely under both histories without distinguishing them, and no
admissibility restriction is needed. What overflow costs is availability, not confidentiality: a
padded refusal discards the caller's write. \Cref{prop:capacity} states this as the
capacity--confidentiality trade-off --- padded refusal (chosen here) hides the load at the price of
lost writes, silent refusal keeps the writes at the price of the load --- and \cref{sec:eval-refusal}
measures both, including the utilization at which the losses begin.
\end{enumerate}

Every component is thus a function of the public parameters, the leakage of the operations issued,
and randomness independent of $b$. A simulator given only $\mathcal{L}$ reproduces the view: it
samples the relabeling $\pi$ itself, emits fixed-length uniform strings for the snapshot, and for
each operation emits the region indices $\mathcal{L}^{\mathrm{op}}$ names under $\pi$, drawing the
uniform padding and per-call probe randomness itself. Hence the real and simulated views are
indistinguishable and $\Pr[b'=b] = \tfrac{1}{2}$ in $\mathrm{G}_2$, giving
\[
\AdvH_{\Pi,\mathcal{A}}(\lambda) \le \mathsf{Adv}^{\mathrm{prf}}_{F}(\lambda) +
\mathsf{Adv}^{\mathrm{ind\$}}_{\mathsf{AE}}(\lambda) \le \negl(\lambda).
\]
\hfill$\blacksquare$

\paragraph{Predicate hiding (\cref{cor:predicate})} Let $\mathcal{D}$ be a distribution over
cohort-H histories, let $\phi$ be a predicate on cohort-H contents, and let $\mathcal{P}$ be the
best predictor of $\phi$ that does not see the adversary's view. Suppose an adversary $\mathcal{B}$
predicts $\phi$ from its view with advantage $\delta$ over $\mathcal{P}$, where the advantage is
taken in expectation over $\mathcal{D}$.

Construct a distinguisher $\mathcal{A}$ for $\mathbf{G}^{\mathrm{H}}$ as follows. $\mathcal{A}$
samples $\sigma_0 \leftarrow \mathcal{D} \mid \phi = 0$ and $\sigma_1 \leftarrow \mathcal{D} \mid
\phi = 1$, submits them, runs $\mathcal{B}$ on the resulting view, and outputs $\mathcal{B}$'s
prediction of $\phi$ as its guess for $b$. Conditioned on the sampled pair, $\mathcal{A}$ is correct
exactly when $\mathcal{B}$ predicts $\phi$ correctly, so
$\Pr[b' = b] = \tfrac{1}{2} + \tfrac{1}{2}\,\delta_{\sigma_0,\sigma_1}$ where
$\delta_{\sigma_0,\sigma_1}$ is $\mathcal{B}$'s advantage on that pair. Averaging over the draw of
the pair gives $\AdvH_{\Pi,\mathcal{A}} \ge \delta/2$, hence
$\delta \le 2\,\AdvH_{\Pi,\mathcal{A}}(\lambda)$. The averaging is necessary because $\mathcal{B}$'s
advantage may be concentrated on a small set of histories; the bound holds for the average, not
pointwise. This is the standard equivalence between semantic security and indistinguishability,
instantiated for the cohort-H game.

The direction of the corollary is worth restating, because it is easy to invert by accident. It
converts a \emph{proof} of dual-epoch security into a bound on predicate recovery. It does not
convert a \emph{measurement} into such a bound: an experiment exhibits one attacker, and the
advantage it achieves is a lower bound on the leakage present, never an upper bound on the leakage
available.

\paragraph{Necessity of the conditions} Condition S2 requires pseudorandom rather than merely
CPA-secure ciphertexts because free slots are uniform strings that must be indistinguishable from
live ones. Condition S3's whole-region rewrite is used in item 3: a per-slot write would reveal
which slot changed and, accumulated, the occupancy that item 2 hides. S3's padded refusal is used in
item 6 and is what lets the theorem cover overflow without an admissibility assumption; dropping it
for silent refusal reopens the capacity channel of \cref{prop:capacity}, which
\cref{sec:eval-refusal} measures. The flush at $\tau$ required by the game closes the
staged-remainder channel. \Cref{prop:seal,prop:fp} show that dropping S1, S2 or S3 admits a
distinguisher, which \cref{sec:eval-game} realizes as ablations, together with the observation that
the keyed relabeling is \emph{not} among the conditions and that removing it leaves the game at
chance.

\section{Proof of Theorem~\ref{thm:reactive}}
\label{supp:reactive}

We prove that a store satisfying S1, S2 and the clock conditions C1--C3 is reactively
$\mathcal{L}_{\mathrm{clock}}$-dual-epoch secure (\cref{def:reactive}) against every polynomial-time
feedback agent $\mathsf{P}$. The structure mirrors the proof of \cref{thm:main}; what changes is that
the operation sequence after $\tau$ is produced by $\mathsf{P}$ reacting to plaintext search results,
so the argument must show that this reaction never reaches the adversary's view.

\paragraph{Setup} Fix $\mathsf{P}$ and an adversary $\mathcal{A}$. In
$\mathbf{G}^{\mathrm{R}}_{\Pi,\mathsf{P},\mathcal{A}}$ the challenger builds cohort~H from $\sigma_b$,
flushes at $\tau$, and then for each adversary-supplied input runs one turn of $\mathsf{P}$: a
sequence of $\mathsf{search}$ operations whose plaintext results $\mathsf{P}$ reads, followed by
$\mathsf{put}$, $\mathsf{delete}$ or $\mathsf{update}$ operations that may depend on those results.
The adversary receives $\mathcal{L}^{\mathrm{op}}$ after each turn and, separately, the clock leakage
$\mathcal{L}_{\mathrm{clock}}$ of unlock, lock and flush times.

\paragraph{Game $\mathrm{H}_0$} The real reactive game.

\paragraph{Game $\mathrm{H}_1$} Replace every pseudorandom function under a store key by a lazily
sampled random function, as in $\mathrm{G}_1$. The gap is $\mathsf{Adv}^{\mathrm{prf}}_F(\lambda)$ up
to the number of derived keys. This step is oblivious to who chose the operations, so $\mathsf{P}$'s
reactivity does not affect it.

\paragraph{Game $\mathrm{H}_2$} Replace every ciphertext written to disk by a uniformly random string
of the same length (C2), keeping plaintexts in challenger state so no stored ciphertext is decrypted.
The gap is $\mathsf{Adv}^{\mathrm{ind\$}}_{\mathsf{AE}}(\lambda)$. After this step every byte on disk,
in every slot and in the sealed journal that carries staged records between flushes, is a uniform
string whose length is a public parameter.

\paragraph{The view is a function of the clock} We show the adversary's view in $\mathrm{H}_2$ is
determined by $\mathcal{L}_{\mathrm{clock}}$ and public randomness, independently of $b$, of the
inputs, and of every branch $\mathsf{P}$ took.

\begin{enumerate}
\item \emph{Snapshot.} By S1 and $\mathrm{H}_2$ the snapshot is a tuple of fixed-length uniform
strings whose count and lengths are public, identical under $b=0$ and $b=1$.
\item \emph{Writes.} By C1 the set of disk locations a flush rewrites is the image of the flush index
under the deterministic sweep: the $\kappa$ slots after a cursor that advances by $\kappa$ each flush from a
start fixed at unlock. This set is a function of the flush index alone. It does not depend on how many
records $\mathsf{P}$ staged, on their content, or on whether $\mathsf{P}$ wrote at all --- a flush
occurs every $\Delta$ regardless (C1). By $\mathrm{H}_2$ the bytes written are uniform. So the write
transcript is the clock's sequence of cursor ranges, filled with uniform bytes.
\item \emph{Reads.} By C3 a $\mathsf{search}$ issues no disk access after the unlock scan. Therefore
the $\mathsf{search}$ operations $\mathsf{P}$ performs, and the results it reads, produce \emph{no}
storage event. This is the crux: the only place $\mathsf{P}$'s secret-dependent behaviour could enter
the transcript is through the operations it issues, and its reads issue nothing while its writes are
placed by the clock, not by content. The branch $\mathsf{P}$ takes is thus invisible.
\item \emph{The journal.} Staged records wait in the fixed-size sealed journal, which C1 rewrites in
full every flush whether it holds staged records or none; by $\mathrm{H}_2$ its bytes are uniform of
fixed length. Its rewrites are on the clock and carry nothing.
\item \emph{Deletions and updates.} A deletion writes nothing at call time (C1 admits only the
clocked sweep as a writer); its slot is freed when the sweep reaches it, an event already accounted
in item~2. An update is a deletion and a staged write, both covered.
\end{enumerate}

Every component is a function of the unlock/lock/flush clock and uniform bytes. A simulator given
$\mathcal{L}_{\mathrm{clock}}$ emits the fixed-length uniform snapshot, and at each flush time emits
the cursor range that flush index determines, filled with uniform bytes; it runs neither $\mathsf{P}$
nor the challenge, and never sees $b$ or the inputs. Hence the real and simulated views are
indistinguishable and
\[
\AdvR_{\Pi,\mathsf{P},\mathcal{A}}(\lambda) \le \mathsf{Adv}^{\mathrm{prf}}_F(\lambda) +
\mathsf{Adv}^{\mathrm{ind\$}}_{\mathsf{AE}}(\lambda) \le \negl(\lambda).
\]
\hfill$\blacksquare$

\paragraph{The lock} C1 includes one flush at lock, always. A lock is a time the adversary chose
(\cref{def:reactive}), so its flush is a clock event like any other, and the cursor range it writes
is the next one in the sweep. A store that flushed \emph{more} than once at lock when more records
were staged than its journal holds would make the number of lock-time flushes a function of how much
$\mathsf{P}$ wrote, which violates C1; this is why Mnemos-M's back-pressure, which refuses a
$\mathsf{put}$ to the caller rather than stage beyond the journal, is the configuration
\cref{cor:m-reactive} covers without exception.

\paragraph{Part (i)} Let $\Pi$ be operation-oblivious and let $\mathsf{P}$ issue at most $q$ turns.
Define hybrids $\mathrm{T}_0,\dots,\mathrm{T}_q$, where in $\mathrm{T}_j$ the operations of the
first $j$ turns are replaced by the empty sequence at the same times and the remaining turns run
$\mathsf{P}$ as in the real game. Consecutive hybrids differ in the operations of one turn, and by
obliviousness the adversary's view of that turn, given everything before it, is indistinguishable
under the two; the later turns are a function of that view and of $\mathsf{P}$'s state, which the
adversary never sees. So $|\Pr[\mathrm{T}_{j-1}\Rightarrow 1] - \Pr[\mathrm{T}_j\Rightarrow 1]| \le
\negl(\lambda)$, and in $\mathrm{T}_q$ the view is the snapshot followed by an idle schedule, both
independent of $b$ by S1, S2 and the obliviousness of idle turns across histories. Summing, $\AdvR
\le q\cdot\negl(\lambda)$.

\paragraph{Part (ii)} Fix a store state and two operation sequences $O$, $O'$, searches included,
whose leakage is at statistical distance $\varepsilon$, and let $T$ be the optimal test between the
two leakage distributions, with advantage $\varepsilon$. Let $x^\star$ be a fact absent from that
state, let $\sigma_1$ produce the state and $\sigma_0$ produce it with $x^\star$ added, both padded to
equal length. $\mathsf{P}$, on input $x$, searches for $x$ and issues $O$ if the top result is a
near-duplicate of $x$ and $O'$ otherwise. $\mathcal{A}$ supplies $x^\star$ and outputs $b'=0$ exactly
when $T$ says $O$. Under $b=1$ no near-duplicate exists and $\mathsf{P}$ issues $O'$; under $b=0$ it
issues $O$ except with the probability $\delta$ that the search misses the stored duplicate. Two
further differences remain between the worlds: the leakage of $\mathsf{P}$'s search, and the
difference between the two store states on which $O$ and $O'$ run. An adversary of
\cref{def:dualepoch} can issue the same search and the same operations on either state itself, so by
dual-epoch security both differences are indistinguishable. Hence $\Pr[b'=b] \ge \tfrac12 +
\tfrac{\varepsilon}{2} - \delta - \negl(\lambda)$, and $\AdvR \ge \varepsilon/2 - \delta -
\negl(\lambda)$. \Cref{prop:derived} is the case where $O$ is a write and $O'$ is nothing.

\paragraph{Why Mnemos-P is excluded} The proof uses C1 and C3 exactly where \cref{prop:derived}'s
attack strikes. Mnemos-P flushes on demand, violating C1: a flush occurs \emph{because} $\mathsf{P}$
staged a record, so the occurrence of a flush is the write-occurrence bit a deduplicating agent
leaks. And Mnemos-P places by content, so even under a clocked flush the region a reactive write
lands in would depend on the retrieved plaintext. The clock-only store closes both, which is why the
composition holds for it and not for the partitioned design --- the separation
$\mathrm{RDE}\Rightarrow\mathrm{DE}$, $\mathrm{DE}\not\Rightarrow\mathrm{RDE}$ of
\cref{cor:m-reactive}. \Cref{sec:eval-reactive} measures both sides.


\section{Proofs of Lemma~\ref{lem:slot} and the Separations}
\label{supp:separations}

\paragraph{\Cref{lem:slot}} The placement rule is invariant under any permutation of the slot labels
of $r$, and by hypothesis nothing else distinguishes them, so the labels are exchangeable. Hence the
set $O_b(r)$ of slots cohort H occupies at $\tau$ is a uniformly random subset of size $L_b(r)$,
whatever the history. Post-$\tau$ records draw uniformly without replacement from its complement, so
$k$ of them occupy a uniformly random ordered $k$-subset of the complement of a uniformly random
$L_b(r)$-subset, which is a uniformly random ordered $k$-subset of all $C$ labels whenever
$k \le C - L_b(r)$. The hypothesis that no partition is overfilled guarantees that bound under both
histories, so the distribution does not depend on $b$. By S2 the snapshot does not reveal $O_b(r)$,
so the adversary cannot condition on it.

\paragraph{\Cref{prop:fp}} Take any forward- and Type-II-backward-private dynamic SSE scheme and
adjoin to its client state a per-partition update counter, incremented on each $\mathsf{put}$ and
stored under a pseudorandom handle in the snapshot. Forward and backward privacy are transcript
notions and are unaffected by client-side state, so the augmented scheme keeps them. The snapshot now
reveals a \emph{multiset} of counter values: $m$ records written to one partition yield $\{m\}$,
$m/2$ written to each of two yield $\{m/2,m/2\}$, and no relabeling of partitions maps one to the
other. The distinguisher submits those two histories and reads the multiset off the snapshot, so
$\AdvH = 1/2$.

\paragraph{\Cref{prop:breach}} Breach resistance asks that a single snapshot be simulatable from
$\mathcal{L}^{\mathrm{snap}}$. That is the $\tau$-only restriction of \cref{def:dualepoch}, so a
dual-epoch simulator stopped at $\tau$ is a breach-resistance simulator. For the converse, the
witness of \cref{prop:fp} is breach-resistant with respect to a snapshot leakage that includes the
counter multiset, and forward/backward private, and is not dual-epoch secure with respect to the
same leakage once the post-$\tau$ writes are read jointly with it: induced puts tell the adversary
which counter is which topic, and the multiset then gives cohort H's per-topic load. Neither of the
two notions constrains that joint object.

\paragraph{\Cref{prop:hi}} History independence fixes the distribution of the stored state given its
\emph{contents}, so two insertion orders of one set are indistinguishable. The two challenge
histories of \cref{def:dualepoch} may hold different sets. A content-placed, strongly
history-independent store still holds a topic's records in that topic's region, and a growing
layout reveals each region's size; with $\sigma_0$ placing $m$ records on one topic and $\sigma_1$
spreading them over two, the snapshot's region sizes differ, and the distinguisher wins with
advantage $1/2$. \Cref{prop:seal} isolates the two properties that remove the channel.

\paragraph{\Cref{prop:seal}} With growing regions, a region's size is its occupancy, so the histories
above are distinguished from one snapshot even though no order inside a region is revealed. With
fixed-size regions whose free slots are not sealed, free slots are distinguishable from ciphertext
(for instance, zero bytes), so the boundary between written and unwritten slots again gives each
region's occupancy. \Cref{sec:eval-game} realizes both stores as ablations.

\section{The Counting Study in Full}
\label{supp:counting}

\Cref{prop:fp} separates forward and backward privacy from dual-epoch security by exhibiting state
a forward-private scheme persists locally: a per-keyword update counter, which distinguishes two
histories with different count \emph{profiles}. The main paper reports the headline, that on a
realistic corpus the channel is weak. This section gives the method behind that number.

\paragraph{Setup} The corpus is D3, whose keyword tags follow a heavy-tailed distribution with many
ties, which is the property the attack turns on. Each record carries three tags drawn from a
503-term vocabulary. A Sophos-style index is built over it, and the snapshot exposes, for each
keyword, a pseudorandom handle and the number of updates recorded under it. The adversary never
sees the keyword itself.

\paragraph{The adversary} It holds auxiliary knowledge of keyword frequencies: a fraction $\alpha$
of the corpus in the clear, from which it estimates an expected count for every keyword. It then
matches observed counters to expected ones, greedily and in descending order of expected count,
which is the strongest simple matching available to it. The sweep runs
$\alpha \in \{0.10, 0.25, 0.50\}$ against a range of tag skews.

\paragraph{Why it is weak here} Three effects compound. Most keywords are rare, so most counters are
small integers. Small integers collide: the number of keywords sharing a given count grows sharply
as the count falls, and keywords with equal counts are information-theoretically indistinguishable
from counts alone. And the record-level question the paper actually cares about, the topic of a
stored record, requires recovering several of a record's tags rather than one, so per-keyword errors
compound rather than average out.

The \emph{tie ceiling}, the fraction of keywords whose count is unique, bounds the recovery rate of
any matching based on counts alone. Averaged over the sweep, keyword recovery is 0.05 and
record-level topic accuracy never exceeds the most-frequent-class baseline. The worst case remains
real, which is what makes the proposition a separation rather than an observation, but a deployment
whose client state is sealed and padded to a fixed size removes the channel outright.

\section{Query Equality in Full}
\label{supp:query}

\Cref{sec:eval-rw} reports that neither probe mechanism hides query equality from an adversary that
thresholds read-set overlap. This section gives the measurement in full (\cref{tab:supp-e15}), the
reason cover does not help, and the anonymity-set analysis, which is the property that separates
the two mechanisms.

\begin{table*}[t]
\centering
\caption{Query equality against three adversaries ($B=64$, $k_{\mathrm{probe}}=4$), over 20 independent stores with 16 queries each, every query searched 10 times in random order. The overlap threshold $\tau$ and the clustering threshold are fixed on pairs among half of each store's queries and scored on the other half, so no adversary is fitted to the pairs it is scored on. $J$ is Youden's (TPR $-$ FPR); ARI is the adjusted Rand index of the recovered query groups. Applies to the RW class only.}
\label{tab:supp-e15}
\small
\setlength{\tabcolsep}{4pt}
\begin{tabular}{lrrrrr}
\toprule
Probe mechanism & Regions/search & Exact-equality $J$ & Threshold $J$ & $\tau$ (mode) & Clustering ARI \\
\midrule
Genuine probes only ($s=1$) & 4 & 0.99\,[0.99, 1.00] & 0.98\,[0.97, 1.00] & 4 & 0.94\,[0.89, 0.99] \\
Randomised superset ($s=2$, shipped) & 8 & 0.00\,[0.00, 0.00] & 0.98\,[0.97, 0.99] & 4 & 0.66\,[0.59, 0.71] \\
Randomised superset ($s=4$) & 16 & 0.00\,[0.00, 0.00] & 0.59\,[0.57, 0.60] & 6 & 0.13\,[0.11, 0.15] \\
Grouped probes (group size 8) & 8 & 0.85\,[0.82, 0.88] & 0.85\,[0.82, 0.88] & 1 & 0.59\,[0.53, 0.65] \\
\bottomrule
\end{tabular}
\end{table*}


\paragraph{The adversary} It observes, for each search, the set of regions read, and must decide
whether two searches carried the same query. It never sees a query or a result. The weak test is
exact equality of the two read sets; the strong one compares the size of their intersection,
$|R_1 \cap R_2|$, against a threshold $\tau$; the third clusters the graph that joins searches
whose overlap reaches $\tau$. Each threshold is fixed on half of a store's queries and scored on
the other half, over \REfifteenRepeats{} independent stores, so no adversary is credited with a
threshold fitted to the answers.

\paragraph{The randomized superset} A search reads $s\,k$ regions ($k = k_{\mathrm{probe}}$), of
which $k$ are the genuine nearest partitions of the query and the remaining $(s-1)k$ are drawn with
fresh randomness on every call. Two searches on the same query therefore almost never produce
identical read sets, and the exact-equality test scores $J = \REfifteenExactStwo$ at $s = 2$,
which is what makes it look as though this mechanism closes the channel.

It does not, because the genuine probes survive every repeat. Treating cover as uniform, two
searches for the same query share the $k$ genuine regions plus about $(s-1)^2k^2/B$ cover
collisions, and two unrelated queries share about $s^2k^2/B$, so the expected gap is
\[
  k\Bigl(1 - \frac{(2s-1)\,k}{B}\Bigr).
\]
More cover narrows the gap linearly in $s$ while the variance of both overlaps grows, and the gap
closes only when $s\,k$ approaches $B/2$, where every search reads about half the store. At the
shipped parameters the threshold test separates the populations with $J = \REfifteenThreshStwo$;
doubling the cover to $s = 4$ lowers it to $J = \REfifteenThreshSfour$. We had registered the
expectation that more cover would barely matter, with a lower bound of 0.60 on this interval, and
it does matter: cover is a real knob. It is a costly one, though, since read volume grows in step
with $s$ and the channel stays open until a search reads half the store, at which point reading
all of it, obliviously, costs little more.

\paragraph{The grouped mechanism} The partitions are divided into fixed disjoint groups at setup,
and a search reads exactly one whole group. Two searches therefore read either the same regions or
disjoint ones, so repeated-query recognition succeeds at the granularity of a group and
$J = \REfifteenThreshGrouped$: every repeat is recognised, and two distinct queries are confused
exactly when they fall in one group.
What it buys in exchange is a guarantee the randomized superset cannot state: within a group every
query is indistinguishable from every other, giving a deterministic anonymity set equal to the
group size, independent of the adversary's effort. The randomized superset's protection is
statistical and degrades with the number of observed repeats; the grouped mechanism's is structural
and does not. Together with the permutation recovery of \cref{sec:eval-permutation}, which the
grouped mechanism defeats outright, this is why it is the one preferred against a read observer
despite recognizing repeats more often.

\section{From Attacker Performance to Information Lower Bounds}
\label{supp:fano}

The leakage ledger converts each channel's measured attacker performance into a lower bound, in
bits, on the mutual information $I(X;V)$ between the secret $X$ the channel is about and the
adversary's view $V$. The conversion is Fano's inequality, and this section states it exactly for
each metric family. In particular an AUC, which is a ranking statistic and not a decoding accuracy,
is never fed to Fano as if it were one.

\paragraph{The inequality} For a secret $X$ over $K$ values and any decoder that recovers it from
$V$ with error probability $P_e$, Fano's inequality gives
$H(X\mid V) \le h(P_e) + P_e\log_2(K-1)$, where $h$ is the binary entropy. Hence
\begin{align*}
I(X;V) &= H(X) - H(X\mid V)\\
       &\ge H(X) - h(P_e) - P_e\log_2(K-1).
\end{align*}
The ledger evaluates this at the \emph{empirical} accuracy $1-P_e$ of the attacker the experiment
exhibited. It is a lower bound twice over: Fano bounds the information any decoder needs, and the
experiment's decoder is one particular decoder, so a better one can only raise $I(X;V)$.

\paragraph{The prior term $H(X)$} The submitted ledger used $H(X)=\log_2 K$. That is exact only for a
\emph{uniform} prior; under a skewed prior $H(X)<\log_2 K$ and the $\log_2 K$ form can exceed the
true $I(X;V)$, so it is not a valid lower bound. The ledger now records, per row, whether the prior
is uniform \emph{by construction} or empirical:
\begin{itemize}
\item \emph{Binary games} (E5, the refusal channel E9, query equality E15, dedup E3): $K=2$ and the
secret is a fair challenge bit or a balanced held-out label, so $H(X)=1$ exactly.
\item \emph{Permutation recovery} (E7): the secret is a region's logical partition under a
\emph{keyed} permutation, uniform over the $B$ regions by construction, so $H(X)=\log_2 B$ exactly.
This is why the headline ``\RLedgerPermBits{} of \RLedgerPermBitsMax{} bits of partition identity''
is a valid bound rather than an artifact of assuming uniformity.
\item \emph{Topic and region recovery} (E4, E6, E8): the prior is the corpus's topic (or region)
distribution, which is skewed. The ledger uses the empirical entropy of that distribution, recorded
by the experiment; where a run predates this recording the row is flagged
\texttt{uniform (assumed)} and its bound is the conservative $\log_2 K$ value, which for these rows
overstates rather than understates and is marked as such.
\end{itemize}

\paragraph{Metric families, and why AUC is not converted}
\begin{itemize}
\item \emph{Balanced accuracy} (binary games, dedup): used directly as $1-P_e$, with the threshold
fixed on a training half and the accuracy measured on a held-out half, so it is the accuracy the
attacker \emph{commits} to, not an in-sample optimum.
\item \emph{Youden's $J$} (refusal, query equality): for a single balanced pair the balanced
accuracy of the thresholded test is exactly $(1+J)/2$, so $J$ maps to accuracy without a ranking
assumption.
\item \emph{Advantage over chance} (topic recovery): the accuracy is $\text{chance}+\text{advantage}$
directly.
\item \emph{Precision} (multi-version): the probability the named region is correct is used as the
decoding accuracy for a $K=B$ secret.
\item \emph{AUC} is \textbf{not} converted. AUC measures ranking, not decoding, and the earlier
$1/4+\text{AUC}/2$ bound credited a perfect classifier with only $0.75$ accuracy while crediting a
useless one with more than chance --- it is loose in one direction and unfounded in the other. Where
a channel is measured only by AUC and no held-out accuracy was recorded, the ledger reports the AUC
and \emph{no} bit, rather than a fabricated one. Every headline bit in the paper comes from a row
with a recorded accuracy or a $J$.
\end{itemize}

\paragraph{Units} The ledger states, per row, whether the secret is a binary history bit, a topic
label, a region index, or a pair-equality bit, and $\log_2 K$ (or the empirical $H(X)$) is the
ceiling against which the bits are read. A row's bound is never compared across units without that
ceiling stated beside it.

\section{Replication Units and Attacker Specification}
\label{supp:units}

\Cref{tab:supp-units} names, for each experiment, the unit an interval resamples and how many
independent units there are. An interval is only as wide as the unit is independent: resampling
records within one store measures sampling noise in that store, not the variation between stores,
so every row resamples the outermost unit that was drawn afresh. Means over at least 20 units carry
a percentile bootstrap interval (10{,}000 resamples); means over fewer carry a Student-$t$ interval,
because a percentile bootstrap over three to five values has only a handful of distinct resample
means and undercovers. AUC intervals resample the held-out units. Every threshold, orientation or
classifier is fitted on a training half of the units and scored on the other half; no parameter
is tuned on the units it is scored on.

\begin{table}[t]
\centering
\caption{The independent unit behind every interval. ``Fresh'' means the unit draws its own
store, cohort split and attacker fit.}
\label{tab:supp-units}
\footnotesize
\setlength{\tabcolsep}{3pt}
\begin{tabular}{@{}llrl@{}}
\toprule
Exp. & Unit (fresh per unit) & $n$ & Interval \\
\midrule
E1 & trial (store, split) & 20 & bootstrap \\
E2 & induction sequence (split, order, fit) & 20 & bootstrap \\
E3 & probe & 80/class & bootstrap, AUC \\
E4 & centroid seed (table, store) & 5/gap & Student-$t$ \\
E5 & game instance & 1{,}000/bit & bootstrap, AUC \\
E6 & seed (store build) & 3 & Student-$t$ \\
E7 & store & 20 & bootstrap \\
E8 & trial (store, versions) & 20 & bootstrap \\
E9 & trial (store) & 20 & bootstrap \\
E15 & store (query split) & 20 & bootstrap \\
E16 & seed (PostgreSQL cluster) & 3 & Student-$t$ \\
E17 & trial (store, versions) & 20 & bootstrap \\
E18 & probe & 30/class & bootstrap, AUC \\
E20 & probe & 400 & bootstrap, AUC \\
E21 & session & 40 & descriptive \\
E22 & store & 20/point & bootstrap \\
E23 & grid point & 1 & descriptive \\
\bottomrule
\end{tabular}
\end{table}

The two experiments with three units (E6, E16) are the deployed-store measurements, where one unit
is a full store build under a stacked encrypted file system. Their intervals are correspondingly
wide, and a null there is reported as ``not detected'' with its interval, not as a closed channel.

\paragraph{The learned attacker} The dual-epoch game (E5) uses one learned distinguisher, fixed
before any arm was run: scikit-learn's
\texttt{HistGradientBoostingClassifier} with 200 iterations, maximum depth 3 and
\texttt{random\_state} 0, and no other hyperparameter set or searched. Its features are the
per-instance trace summary the arm computes from the observer's view (per-region write counts and
their spread, and for slot-granular arms the number of distinct slots touched and the mean, maximum
and standard deviation of their indices). The instances of both bits are shuffled together and
split in half; the classifier is fitted on one half and its AUC and balanced accuracy, at a
threshold fixed on the training half, are reported on the other. Each arm also reports one named
hand-built statistic, oriented on the training half, so that a weak learner cannot make an arm look
closed without the hand-built statistic disagreeing.

The reactive game (E20) has a one-dimensional observation per probe, so it uses no classifier. It
scores three named trace statistics (flushes per turn, region writes per turn, and writes landing
in the probe's content region), each
oriented on the training half of the probes, and reports the one whose held-out AUC departs
furthest from one half. Taking the maximum of three on the held-out half biases the reported AUC
away from chance, which is the conservative direction for an arm claimed to be at chance: the clock-only
arm is at chance even under that selection.

\section{Full Results}
\label{supp:full}

The main paper tabulates the rows its arguments rest on. The tables here give every row, and are
generated from the same result files by the same code, so the two cannot disagree.

\begin{table*}[t]
\centering
\caption{Every channel in one unit. \emph{Bits} is a Fano lower bound on the mutual information between the secret and what the adversary observes, per observation of the unit the channel is about (a record's topic, a region's partition, one game instance), taken at the lower end of the experiment's 95\% interval so that it holds at the experiment's confidence; $^\ast$ marks a result with no interval, for which only the point bound exists, and $^\dagger$ a channel measured only by AUC, which is a ranking statistic and carries no decoding-accuracy bound (the supplement derives the conversion). A test bounds leakage from below, so for a \emph{not detected} channel the informative number is the last column: the smallest effect, in the experiment's own unit, it had 80\% power to detect. Every measured channel, including every ablation.}
\label{tab:supp-ledger}
\footnotesize
\setlength{\tabcolsep}{4pt}
\begin{tabular}{llllrl}
\toprule
Channel & Design & Class & Status & Bits $\geq$ & Detectable $\geq$ \\
\midrule
structure transfer, historically aware adversary & growing layout & W & not det. & 0.000 & 0.027 (accuracy over majority) \\
structure transfer, naive (current content) adversary & growing layout & W & not det. & 0.000 & 0.004 (accuracy over majority) \\
snapshot + post-$\tau$ writes (the game) & Mnemos-P & W & not det. & 0.002 & 0.052 (AUC) \\
game, S1 removed (growing regions) & Mnemos-P ablation & W & \textbf{leaks} & 1.000 & -- \\
game, S2 removed (unsealed slots) & Mnemos-P ablation & W & \textbf{leaks} & 0.784 & -- \\
game, per-slot writes, random slot & Mnemos-P ablation & W & not det. & 0.006 & 0.051 (AUC) \\
game, no keyed permutation & Mnemos-P ablation & W & not det. & 0.000 & 0.052 (AUC) \\
refusal, padded & Mnemos-P & W & not det. & 0.001 & 0.051 (AUC) \\
refusal, silent & Mnemos-P ablation & W & \textbf{leaks} & 0.682 & -- \\
snapshot + post-$\tau$ writes (the game) & Mnemos-M & W & not det. & 0.001 & 0.051 (AUC) \\
game, content placement & Mnemos-M ablation & W & not det. & 0.008 & 0.051 (AUC) \\
game, first-free slot, whole-region writes & Mnemos-P ablation & W & not det. & 0.000 & 0.052 (AUC) \\
game, per-slot writes, first-free slot & Mnemos-P ablation & W & \textbf{leaks} & 1.000 & -- \\
structure transfer, deployed store: append\_control & append\_control & W & not det. & 0.000 & 0.002 (advantage) \\
structure transfer, deployed store: chroma & chroma & W & not det. & 0.004 & 0.020 (advantage) \\
structure transfer, deployed store: faiss & faiss & W & not det. & 0.000 & undetermined \\
structure transfer, deployed store: lancedb & lancedb & W & not det. & 0.000 & undetermined \\
structure transfer, deployed store: lancedb\_indexed & lancedb\_indexed & W & not det. & 0.000 & undetermined \\
structure transfer, deployed store: sqlitevec & sqlitevec & W & not det. & 0.000 & 0.005 (advantage) \\
permutation recovery (spread\_s1) & Mnemos-P & RW & \textbf{leaks} & 4.905 & -- \\
permutation recovery (spread\_s2) & Mnemos-P & RW & \textbf{leaks} & 1.827 & -- \\
permutation recovery (grouped) & Mnemos-P & RW & \textbf{leaks} & 0.000 & -- \\
multi-version diff, g=1 & Mnemos-P & W & \textbf{leaks} & 0.524 & -- \\
multi-version diff, g=2 & Mnemos-P & W & \textbf{leaks} & 0.173 & -- \\
multi-version diff, g=4 & Mnemos-P & W & \textbf{leaks} & 0.049 & -- \\
multi-version diff, g=16 & Mnemos-P & W & \textbf{leaks} & 0.002 & -- \\
multi-version diff, g=64 & Mnemos-P & W & \textbf{leaks} & 0.000 & -- \\
multi-version diff, across a session, rewrite\_at\_lock off & Mnemos-P & W & \textbf{leaks} & 0.031 & -- \\
multi-version diff, across a session, rewrite\_at\_lock on & Mnemos-P & W & not det. & 0.000 & undetermined \\
refusal, silent & Mnemos-P & W & \textbf{leaks} & 0.763 & -- \\
refusal, padded & Mnemos-P & W & not det. & 0.000 & undetermined \\
query equality (spread\_s1), overlap threshold & Mnemos-P & RW & \textbf{leaks} & 0.887 & -- \\
query equality (spread\_s2), overlap threshold & Mnemos-P & RW & \textbf{leaks} & 0.890 & -- \\
query equality (spread\_s4), overlap threshold & Mnemos-P & RW & \textbf{leaks} & 0.252 & -- \\
query equality (grouped), overlap threshold & Mnemos-P & RW & \textbf{leaks} & 0.566 & -- \\
dedup membership (threshold) & Mnemos-P & W & \textbf{leaks} & 1.000 & -- \\
dedup membership (threshold) & Mnemos-M & W & not det. & 0.000 & undetermined \\
dedup membership (llm) & Mnemos-P & W & \textbf{leaks} & 0.041 & -- \\
dedup membership (llm) & Mnemos-M & W & not det. & 0.000 & undetermined \\
\bottomrule
\end{tabular}
\end{table*}

\begin{table*}[t]
\centering
\caption{Exposure of the pre-compromise cohort on a timestamped stream (2,501 records over 4 months, $\tau$ at month 2, 10 topics, chance 0.100). Source: agent-authored memories. \emph{Membership} is how the attack learns which region holds a given earlier record: the growing layout's snapshot discloses it, while for the sealed designs the harness supplies it, so those rows state what the labels would be worth to an adversary that obtained membership elsewhere (\cref{sec:eval-multisnapshot}). The RW adversary against Mnemos-P recovers the keyed permutation at the rate \cref{sec:eval-permutation} measured, 0.564, and labels partitions from public text only.}
\label{tab:e1}
\small
\setlength{\tabcolsep}{4pt}
\begin{tabular}{llrlr}
\toprule
Design & Class & Budget & Membership & \advhist{} [95\% CI] \\
\midrule
Growing layout & W & 4 & snapshot & 0.055\,[0.046, 0.063] \\
Growing layout & RW & 4 & snapshot & 0.055\,[0.046, 0.063] \\
Mnemos-M & W & 4 & supplied & -0.001\,[-0.003, 0.001] \\
Mnemos-M & RW & 4 & supplied & -0.000\,[-0.002, 0.002] \\
Growing layout & W & 16 & snapshot & 0.064\,[0.056, 0.071] \\
Growing layout & RW & 16 & snapshot & 0.064\,[0.056, 0.071] \\
Mnemos-M & W & 16 & supplied & -0.001\,[-0.003, 0.001] \\
Mnemos-M & RW & 16 & supplied & -0.002\,[-0.004, 0.000] \\
Growing layout & W & 64 & snapshot & 0.097\,[0.092, 0.103] \\
Growing layout & RW & 64 & snapshot & 0.097\,[0.092, 0.103] \\
Mnemos-M & W & 64 & supplied & 0.001\,[0.000, 0.002] \\
Mnemos-M & RW & 64 & supplied & 0.000\,[-0.001, 0.001] \\
Mnemos-P & W & 4 & supplied & 0.011\,[0.006, 0.017] \\
Mnemos-P & RW & 4 & supplied & 0.070\,[0.057, 0.083] \\
Mnemos-RS & W & 4 & supplied & 0.015\,[0.010, 0.020] \\
Mnemos-P & W & 16 & supplied & 0.035\,[0.027, 0.043] \\
Mnemos-P & RW & 16 & supplied & 0.076\,[0.068, 0.084] \\
Mnemos-RS & W & 16 & supplied & 0.039\,[0.033, 0.045] \\
Mnemos-P & W & 64 & supplied & 0.054\,[0.050, 0.059] \\
Mnemos-P & RW & 64 & supplied & 0.097\,[0.085, 0.108] \\
Mnemos-RS & W & 64 & supplied & 0.065\,[0.060, 0.072] \\
\bottomrule
\end{tabular}
\end{table*}

\begin{table*}[t]
\centering
\caption{Every budget measured. The attack on 3,666 memories a local model actually wrote (10 topics, chance 0.100). $w_{\mathrm{real}}$ is the burst size: how many same-topic memories one writing turn contributes to a flush. The growing layout discloses each record's partition in the snapshot and is the positive control; the advantage is scored for an adversary that guesses where it has no label.}
\label{tab:supp-e2}
\small
\setlength{\tabcolsep}{4pt}
\begin{tabular}{lrrr}
\toprule
Layout & $w_{\mathrm{real}}$ & Turns/topic & \advhist{} [95\% CI] \\
\midrule
growing (control) & 1 & 4 & 0.041\,[0.024, 0.058] \\
growing (control) & 1 & 8 & 0.037\,[0.016, 0.062] \\
growing (control) & 1 & 16 & 0.027\,[0.009, 0.047] \\
growing (control) & 1 & 32 & 0.051\,[0.017, 0.091] \\
growing (control) & 2 & 4 & 0.109\,[0.085, 0.130] \\
growing (control) & 2 & 8 & 0.076\,[0.057, 0.094] \\
growing (control) & 2 & 16 & 0.149\,[0.113, 0.179] \\
growing (control) & 2 & 32 & 0.158\,[0.119, 0.191] \\
growing (control) & 4 & 4 & 0.085\,[0.076, 0.094] \\
growing (control) & 4 & 8 & 0.182\,[0.172, 0.193] \\
growing (control) & 4 & 16 & 0.184\,[0.172, 0.196] \\
growing (control) & 4 & 32 & 0.188\,[0.173, 0.203] \\
Mnemos-P & 1 & 4 & 0.003\,[-0.002, 0.010] \\
Mnemos-P & 1 & 8 & 0.004\,[-0.001, 0.010] \\
Mnemos-P & 1 & 16 & 0.004\,[-0.001, 0.009] \\
Mnemos-P & 1 & 32 & 0.009\,[0.000, 0.018] \\
Mnemos-P & 2 & 4 & 0.018\,[0.010, 0.025] \\
Mnemos-P & 2 & 8 & 0.022\,[0.015, 0.030] \\
Mnemos-P & 2 & 16 & 0.031\,[0.020, 0.041] \\
Mnemos-P & 2 & 32 & 0.014\,[0.009, 0.020] \\
Mnemos-P & 4 & 4 & 0.045\,[0.033, 0.058] \\
Mnemos-P & 4 & 8 & 0.065\,[0.054, 0.074] \\
Mnemos-P & 4 & 16 & 0.045\,[0.035, 0.057] \\
Mnemos-P & 4 & 32 & 0.021\,[0.014, 0.028] \\
\bottomrule
\end{tabular}
\end{table*}

\begin{table*}[t]
\centering
\caption{Cost at every size measured. Against the in-tree Rust plaintext IVF, which shares the quantisation and the partitioner, so a ratio is a statement about the security machinery. \emph{Put} is the sparse regime: one record, one flush, which is what an agent writing seconds apart pays. \emph{kB written/record} is what a whole-region rewrite costs per record; a dash marks a design that does not count its bytes or keeps nothing on disk.}
\label{tab:supp-e10}
\footnotesize
\setlength{\tabcolsep}{4pt}
\begin{tabular}{lrrrrrrr}
\toprule
Design & $n$ & Search $p50$ (ms) & Put $p50$ (ms) & kB written/record & RAM (MB) & Disk (MB) & Unlock (s) \\
\midrule
Plaintext IVF (Rust) & 1,000 & 0.091 & 0.04 & -- & 0.0 & 0.0 & 0.00 \\
Mnemos-M, exact & 1,000 & 0.137 & 4.33 & 198 & 5.5 & 5.8 & 0.21 \\
Mnemos-M, in-memory IVF & 1,000 & 0.068 & 3.81 & 198 & 5.5 & 5.8 & 0.20 \\
Mnemos-P & 1,000 & 1.237 & 15.00 & 198 & 0.0 & 5.8 & 0.25 \\
Mnemos-RS & 1,000 & 0.662 & 3.70 & -- & 0.0 & 5.8 & 0.30 \\
Plaintext IVF (Rust) & 3,000 & 0.069 & 0.02 & -- & 0.0 & 0.0 & 0.00 \\
Mnemos-M, exact & 3,000 & 0.454 & 6.04 & 290 & 16.4 & 17.2 & 0.21 \\
Mnemos-M, in-memory IVF & 3,000 & 0.104 & 5.84 & 290 & 16.4 & 17.2 & 0.23 \\
Mnemos-P & 3,000 & 1.718 & 17.24 & 290 & 0.0 & 17.2 & 0.23 \\
Mnemos-RS & 3,000 & 0.823 & 3.92 & -- & 0.0 & 17.2 & 0.32 \\
Plaintext IVF (Rust) & 10,000 & 0.106 & 0.02 & -- & 0.0 & 0.0 & 0.00 \\
Mnemos-M, exact & 10,000 & 1.467 & 8.96 & 478 & 54.5 & 57.1 & 0.26 \\
Mnemos-M, in-memory IVF & 10,000 & 0.238 & 8.12 & 478 & 54.6 & 57.1 & 0.38 \\
Mnemos-P & 10,000 & 2.557 & 22.15 & 478 & 0.0 & 57.1 & 0.28 \\
Mnemos-RS & 10,000 & 1.522 & 5.65 & -- & 0.0 & 57.2 & 0.36 \\
Plaintext IVF (Rust) & 30,000 & 0.288 & 0.05 & -- & 0.0 & 0.0 & 0.00 \\
Mnemos-M, exact & 30,000 & 6.404 & 13.13 & 713 & 163.3 & 170.8 & 0.35 \\
Mnemos-M, in-memory IVF & 30,000 & 0.602 & 14.04 & 713 & 163.5 & 170.8 & 1.21 \\
Mnemos-P & 30,000 & 3.813 & 28.19 & 713 & 0.0 & 170.8 & 0.39 \\
Mnemos-RS & 30,000 & 2.383 & 8.38 & -- & 0.0 & 171.1 & 0.43 \\
Plaintext IVF (Rust) & 100,000 & 0.422 & 0.06 & -- & 0.0 & 0.0 & 0.00 \\
Mnemos-M, exact & 100,000 & 23.194 & 32.18 & 1,185 & 543.7 & 568.2 & 0.79 \\
Mnemos-M, in-memory IVF & 100,000 & 1.778 & 33.85 & 1,185 & 544.3 & 568.2 & 6.67 \\
Mnemos-P & 100,000 & 6.667 & 46.26 & 1,185 & 0.0 & 568.2 & 0.69 \\
Mnemos-RS & 100,000 & 3.442 & 13.09 & -- & 0.0 & 569.4 & 0.71 \\
\bottomrule
\end{tabular}
\end{table*}


\section{Further Tables and Figures}
\label{supp:further}

These results are argued from in the main paper, which carries their headline in prose or in the
leakage ledger (\cref{tab:ledger}); the full tables are here.

\begin{figure*}[t]
\centering
\includegraphics[width=0.78\textwidth]{figures/fig2_threat_model.pdf}
\caption{The $\tau$-arrival adversary. Records written before $\tau$ (cohort H) are available only
as stored state; activity after $\tau$ (cohort P) is observed and partly induced. Where assignment
depends on content, partitions hold records from both cohorts, so labels learned after $\tau$
describe cohort H as well --- but they reach a given cohort-H record only if the layout discloses
which partition holds it, which a growing layout does and a sealed fixed-size one does not
(\cref{lem:transfer}).}
\label{fig:threat}
\end{figure*}

\begin{table}[t]
\centering
\caption{Notation}
\label{tab:notation}
\small
\begin{tabular}{@{}lp{0.68\columnwidth}@{}}
\toprule
Symbol & Meaning \\
\midrule
$B$, $C$ & partitions (or regions); slots per partition \\
$w$ & records written per flush (\emph{epoch\_writes}) \\
$\Delta$, $k$ & clocked flush interval; slots per sweep flush \\
$k_{\mathrm{probe}}$, $\mathit{top\_k}$ & partitions probed per search; records returned \\
$s$ & probe spread \\
$\tau$ & time of compromise \\
H, P & records written before $\tau$; at or after $\tau$ \\
$\beta$ & injection budget: induced writes per topic \\
$g$ & version-capture granularity, in flushes \\
\advhist, \advpost & topic-recovery advantage on cohorts H and P \\
\bottomrule
\end{tabular}
\end{table}

\begin{figure*}[t]
\centering
\includegraphics[width=\linewidth]{figures/fig_crossover.pdf}
\caption{Where the choice between the designs is decided (main paper, \cref{sec:eval-perf}). Three
measures on different scales are shown as small multiples over a shared $n$. Mnemos-M's exhaustive
ranking is linear in $n$ and crosses the partitioned design between \REtenCrossLo{} and
\REtenCrossHi{} records; with its partition index held in memory (hollow markers) it stays ahead
at every size, paid for at unlock, where every record's partition is recomputed. The middle panel,
resident memory, is what bounds Mnemos-M; the index leaves it unchanged.}
\label{fig:crossover}
\end{figure*}

\begin{table*}[t]
\centering
\caption{Advantage on cohort H against the time gap between cohort H and the content the
adversary induces (github, 12 topics, chance
0.083). Gap 0 is an adversary inducing with archival content from cohort H's own
era; the largest gap is one inducing with current content. The i.i.d.\ control is
0.368\,[0.358, 0.379]. Cohort H's topic prior is skewed: its most frequent topic is
0.504 of its records. The last column is accuracy over an adversary that knows
that prior and always guesses the most frequent topic.}
\label{tab:e4}
\small
\setlength{\tabcolsep}{4pt}
\begin{tabular}{llrr}
\toprule
Gap (eras) & Adversary & \advhist{} [95\% CI] & Over majority [95\% CI] \\
\midrule
0 & historically aware & 0.424\,[0.405, 0.443] & 0.004\,[-0.015, 0.023] \\
1 & intermediate & 0.246\,[0.220, 0.275] & -0.174\,[-0.200, -0.145] \\
2 & intermediate & 0.036\,[0.024, 0.046] & -0.384\,[-0.396, -0.374] \\
3 & intermediate & -0.017\,[-0.027, -0.009] & -0.437\,[-0.448, -0.429] \\
4 & intermediate & -0.029\,[-0.034, -0.023] & -0.449\,[-0.455, -0.443] \\
5 & naive (current content) & -0.022\,[-0.026, -0.019] & -0.442\,[-0.446, -0.439] \\
\bottomrule
\end{tabular}
\end{table*}

\begin{table*}[t]
\centering
\caption{Deployed vector stores under a \texttt{gocryptfs} mount, 10
repeats each, 8 topics (chance 0.125). Induced writes are \emph{interleaved across
topics}, so consecutive same-topic writes cannot make recency look like content placement, and
block labels are fitted on one half of them and scored on the other. \advhist{} is the quantity
that matters: whether those labels identify the topics of records written \emph{before} the
compromise.}
\label{tab:e6}
\small
\setlength{\tabcolsep}{4pt}
\begin{tabular}{llrrr}
\toprule
Store & Placed by & Blocks/write & Localisation, held out & \advhist{} \\
\midrule
\texttt{append-control} & time & 2.2 & 0.13\,[0.11, 0.15] & 0.00\,[-0.00, 0.00] \\
\texttt{chroma} & content & 12.5 & 0.14\,[0.11, 0.16] & -0.01\,[-0.02, 0.00] \\
\texttt{faiss} & content & 367.7 & 0.12\,[0.08, 0.17] & 0.00\,[0.00, 0.00] \\
\texttt{lancedb} & time & 12.1 & 0.00\,[0.00, 0.00] & 0.00\,[0.00, 0.00] \\
\texttt{lancedb-indexed} & content & 243.9 & 0.00\,[0.00, 0.00] & 0.00\,[0.00, 0.00] \\
\texttt{sqlitevec} & time & 11.5 & 0.13\,[0.11, 0.15] & -0.00\,[-0.00, 0.00] \\
\bottomrule
\end{tabular}
\end{table*}

\begin{table*}[t]
\centering
\caption{Deduplication membership inference: whether an adversary learns that a fact was already
in cohort H, from the occurrence of a write alone (80 probes per class,
bootstrap intervals). The write rates are the agent's decisions and are the same for both designs;
Mnemos-M's clocked flush makes what the adversary observes identical either way.}
\label{tab:e3}
\small
\setlength{\tabcolsep}{4pt}
\begin{tabular}{llrrr}
\toprule
Decided by & Design & AUC [95\% CI] & Writes, member & non-member \\
\midrule
similarity & Mnemos-P & 1.00\,[1.00, 1.00] & 0.00 & 1.00 \\
similarity & Mnemos-M & 0.50\,[0.50, 0.50] & 0.00 & 1.00 \\
model: qwen2.5 & Mnemos-P & 0.62\,[0.57, 0.67] & 0.74 & 0.97 \\
model: qwen2.5 & Mnemos-M & 0.50\,[0.50, 0.50] & 0.74 & 0.97 \\
\bottomrule
\end{tabular}
\end{table*}

\begin{table*}[t]
\centering
\caption{Recovering the keyed partition permutation from search traffic alone, with no induced
writes, by matching the observed region co-occurrence graph against the one the public table
predicts ($B=64$, $k_{\mathrm{probe}}=4$, 800 queries,
null $1/B = 0.016$), over 20 independent stores.
\emph{FAQ} is the fast approximate QAP solver alone; \emph{refined} adds pairwise-swap local
search from each of its restarts and is the adversary reported, since a leakage measurement is only
as strong as its attacker. Applies to the RW class only.}
\label{tab:e7}
\small
\setlength{\tabcolsep}{4pt}
\begin{tabular}{lrrr}
\toprule
Probe mechanism & Exact, FAQ & Exact, refined & Neighbourhood \\
\midrule
Genuine probes only ($s=1$) & 0.887\,[0.823, 0.934] & 0.963\,[0.897, 1.000] & 0.965 \\
Randomised superset ($s=2$, shipped) & 0.393\,[0.328, 0.460] & 0.564\,[0.469, 0.664] & 1.000 \\
Grouped probes (group size 8) & 0.025\,[0.016, 0.034] & 0.025\,[0.016, 0.034] & 0.172 \\
\bottomrule
\end{tabular}
\end{table*}

\begin{table*}[t]
\centering
\caption{What a sequence of file versions reveals about writes made \emph{before} the
compromise ($B=64$, $w=4$, null precision $1/B=0.0156$).
\emph{Exposure} is how often the record's true region is inside the diff; \emph{precision} is
the chance of naming it by guessing within the diff.}
\label{tab:e8}
\small
\setlength{\tabcolsep}{4pt}
\begin{tabular}{lrrr}
\toprule
Capture & Exposure & Precision [95\% CI] & Candidate regions \\
\midrule
every 1 flushes & 1.000 & 0.208\,[0.207, 0.208] & 4.9 \\
every 2 flushes & 1.000 & 0.108\,[0.108, 0.108] & 9.3 \\
every 4 flushes & 1.000 & 0.058\,[0.058, 0.058] & 17.4 \\
every 16 flushes & 1.000 & 0.022\,[0.022, 0.022] & 46.2 \\
every 64 flushes & 1.000 & 0.016\,[0.016, 0.016] & 63.8 \\
across a session, without re-encryption & 1.000 & 0.0492\,[0.0476, 0.0508] & 20.4 \\
across a session, with re-encryption & 1.000 & 0.0156\,[0.0156, 0.0156] & 64.0 \\
\bottomrule
\end{tabular}
\end{table*}

\begin{table*}[t]
\centering
\caption{The refusal channel ($B=64$, $C=32$, $w=4$;
padding alone rewrites a given region on 0.076 of flushes). The adversary
induces writes destined for a region and watches whether that region is rewritten.}
\label{tab:e9}
\small
\setlength{\tabcolsep}{4pt}
\begin{tabular}{lrrrr}
\toprule
Refusal behaviour & AUC [95\% CI] & Rate, full & Rate, open & Gap \\
\midrule
Silent (writes nothing) & 1.000\,[1.000, 1.000] & 0.070 & 1.000 & +0.930 \\
Padded (rewrites the region anyway) & 0.500\,[0.500, 0.500] & 1.000 & 1.000 & +0.000 \\
\bottomrule
\end{tabular}
\end{table*}

\begin{table*}[t]
\centering
\caption{Capacity under a skewed write trace: when each design first refuses. The partitioned
design leaves the admissible regime of \cref{thm:main} long before the store is full.}
\label{tab:e9cap}
\small
\setlength{\tabcolsep}{4pt}
\begin{tabular}{lrrr}
\toprule
Design & Admitted & First refusal after & Utilisation \\
\midrule
\texttt{p} & 1845 & 488 & 23.8\% \\
\texttt{m} & 2048 & 2048 & 100.0\% \\
\bottomrule
\end{tabular}
\end{table*}

\begin{table*}[t]
\centering
\caption{nDCG@10 against BEIR relevance judgments on the \emph{full} corpora. The columns
separate the two sources of loss: the gap from exhaustive to IVF at the same
$k_{\mathrm{probe}}=4$ is the cost of fixed probing, and the gap from exhaustive to
Mnemos-M, which searches exactly, is the cost of int8 quantisation. \emph{M (IVF)} is Mnemos-M
with its in-memory partition index probing 8 partitions, the number of regions Mnemos-P reads at
$s{=}2$. Refusals during indexing are counted, not swallowed.}
\label{tab:e11}
\small
\setlength{\tabcolsep}{4pt}
\begin{tabular}{lrrrrrrrr}
\toprule
 & & \multicolumn{2}{c}{Plaintext float} & \multicolumn{4}{c}{int8 engines} & \\
\cmidrule(lr){3-4}\cmidrule(lr){5-8}
Domain & $|D|$ & exhaustive & IVF@$k$ & M (exact) & M (IVF) & P ($s{=}2$) & P grouped
 & Refused \\
\midrule
scifact & 5,183 & 0.738 & 0.657 & 0.737 & 0.703 & 0.658 & 0.435 & 0 \\
nfcorpus & 3,633 & 0.370 & 0.313 & 0.369 & 0.342 & 0.319 & 0.256 & 0 \\
fiqa & 57,638 & 0.384 & 0.327 & 0.384 & 0.358 & 0.329 & 0.224 & 0 \\
\bottomrule
\end{tabular}
\end{table*}

\begin{table*}[t]
\centering
\caption{Fetching the $top\_k$ payloads, for the RW class. Path ORAM is a cost model with its
parameters stated ($Z=4$), not a measurement.}
\label{tab:e12}
\small
\setlength{\tabcolsep}{4pt}
\begin{tabular}{lrlrrl}
\toprule
Method & $n$ & Source & $p50$ (ms) & kB read/search & Leaks positions \\
\midrule
resident & 1000 & measured & 0.159 & 0.0 & no \\
position & 1000 & measured & 0.226 & 18.9 & yes \\
read all & 1000 & measured & 3.071 & 5691.1 & no \\
path oram & 1000 & computed & --- & 1665.0 & no \\
resident & 10000 & measured & 1.751 & 0.0 & no \\
position & 10000 & measured & 1.787 & 18.9 & yes \\
read all & 10000 & measured & 30.477 & 56911.4 & no \\
path oram & 10000 & computed & --- & 2119.0 & no \\
resident & 50000 & measured & 12.532 & 0.0 & no \\
position & 50000 & measured & 12.214 & 18.9 & yes \\
read all & 50000 & measured & 154.072 & 283830.3 & no \\
path oram & 50000 & computed & --- & 2573.1 & no \\
\bottomrule
\end{tabular}
\end{table*}


\paragraph{The operating-system artifact audit} \Cref{sec:eval-artifacts} summarises it. On the
evaluation machine the audit recovers change-journal (USN) records naming the store's own files,
confirming that the file system retains pre-$\tau$ write history with timestamps; it enumerates
volume shadow copies and reports where elevation was required to do so; it identifies which
synchronization clients are installed and whether the store lies inside a synchronized tree; and it
reports the hibernation file and pagefile, which matter because a hibernation image captures locked
pages and therefore the unwrapped master key. It also verifies the mechanism a power-event handler
would invoke: after locking, both sealed engines report themselves locked and refuse to search.

\section{Agent-Memory Frameworks, Read from Source}
\label{supp:census}

\Cref{sec:attack-derived} states that four of nine widely used open-source agent-memory
frameworks make whether a write reaches storage depend on whether the content is already stored,
which is the precondition of \cref{prop:derived}. \Cref{tab:supp-census} gives the evidence. The
method is deliberately narrow: each framework's wheel was downloaded from PyPI without its
dependencies and its shipped Python source read directly, and a mechanism is recorded only if it
was located at the cited file and line. Documentation was not used, because documentation
describes intent and the attack depends on behaviour.

The census also corrects a description of mem0 that is widely repeated. mem0 is commonly described as running an update stage that returns
\textsc{noop} for a fact it already holds. The current release has no such stage: it passes the
most similar existing memories to a single extraction call, and then drops any extracted fact
whose hash is already present. The observable consequence is the same --- a known fact never
reaches the vector store --- but it now arises from two mechanisms rather than one, and only the
vector store's files are written conditionally: mem0's message log is written either way. An
adversary that watches the whole store directory sees the difference as \emph{which} files
changed.

\begin{table*}[t]
\centering
\caption{Whether a write reaches storage depends on whether the content is already stored, read from each framework's shipped source (\texttt{pip download --no-deps}, 2026-09-27). \emph{Yes} is recorded only where the mechanism was located at the cited file; anything not located is \emph{none found} or \emph{not verified}, never inferred from documentation. Where it is yes, the source satisfies the precondition of \cref{prop:derived}. Whether the oracle opens when the pipeline runs depends on the deciding model: run end to end it is open in mem0 with \texttt{qwen2.5} and closed in LangMem, which re-stores the memories it retrieves on every call (\cref{tab:frameworks}).}
\label{tab:supp-census}
\small
\setlength{\tabcolsep}{4pt}
\begin{tabular}{llp{3.6cm}lp{4.9cm}}
\toprule
Framework & Version & Default store & Write depends on existing content & Evidence \\
\midrule
mem0 & 2.2.1 & Qdrant, local mode at /tmp/qdrant & \textbf{yes} & \nolinkurl{mem0/memory/main.py:925-955} \\
CrewAI & 1.15.22 & LanceDB & \textbf{yes} & \nolinkurl{crewai/memory/unified_memory.py:93} \\
LangMem & 0.0.30 & LangGraph BaseStore (backend chosen by the caller) & \textbf{yes} & \nolinkurl{langmem/knowledge/extraction.py:195-211} \\
Graphiti (Zep) & 0.30.2 & graph database (caller-configured) & \textbf{yes} & \nolinkurl{graphiti_core/utils/maintenance/combined_extraction.py:171-174} \\
LlamaIndex core & 0.14.25 & caller-chosen vector store & none found & \nolinkurl{llama_index/core/memory/memory_blocks/vector.py:29} \\
Cognee & 1.6.1 & LanceDB (vectors), Kuzu (graph) & not verified & \nolinkurl{cognee/api/v1/config/config.py:215} \\
Agno & 3.0.11 & caller-chosen database & none found & \nolinkurl{agno/db/base.py:580-636} \\
Memobase & 0.0.27 & server-side (client package) & none found & -- \\
Letta & 0.33.2 & server-side & not verified & -- \\
\bottomrule
\end{tabular}
\end{table*}


\section{The Dual-Epoch Game: Per-Ablation Details}
\label{supp:ablations}

\Cref{sec:eval-game} reports the outcome of the game of \cref{def:dualepoch} against the compiled
engines. This section gives the construction of the arms, and why four choices in the experiment are
load-bearing rather than incidental. Each is a way the experiment can appear to measure the
theorem's setting while measuring something else.

\paragraph{Admissibility is verified, not assumed} A history is admissible only if no partition
exceeds its capacity at any point. An experiment that merely \emph{intends} admissibility and wraps
each write in a bare exception handler will silently submit inadmissible histories and measure the
refusal channel, while reporting a number that reads as placement leakage. Every admissible arm here
therefore asserts that the engine's own refusal counter is exactly zero at the end of the arm.

Two distinct defects were caught only by that assertion. The first was an admissibility check that
recomputed each record's partition in floating point by nearest centroid, which disagrees with the
engine's quantised assignment at the margin; the check now calls the engine's own assignment
directly, so the experiment and the implementation cannot differ about where a record went. The
second was subtler: bounding the pre-$\tau$ cohort at the partition capacity left no room for the
adversary's own targeted post-$\tau$ writes, which then overflowed the very partitions the history
had been constructed to keep admissible. The cohort bound is now set below capacity by the
adversary's entire induced budget.

The second defect is the reason the assertion is worth its cost: at 50 instances per bit every
admissible arm reported zero refusals and the experiment looked correct, and only at 1{,}000
instances per bit did an admissible arm report seven. A leakage measurement that is sound at low
power and unsound at high power is indistinguishable from a sound one until the moment it matters.

\paragraph{Evaluation is held out} The distinguisher is fitted on one half of the instances and
scored on the other. Without the split, a classifier with enough features recovers the bit on the
data it was fitted to, and the resulting AUC measures capacity rather than leakage. With it, an AUC
below one half is read as the small-sample anti-learning it is, and is not reported as evidence of a
distinguisher running backwards.

\paragraph{Every arm carries a hand-built statistic} Alongside the gradient-boosted classifier each
arm names one statistic, oriented on the training half, that a reader can check by hand: the spread
of post-$\tau$ writes over regions for the shipped arms, and for the $\lnot$S3 arm the slot-level
features a per-slot write exposes --- the number of distinct slots touched and the mean, maximum and
standard deviation of their indices. A learned AUC with no interpretable companion cannot
distinguish a real channel from a harness artifact; the pair can.

\paragraph{The inadmissible case is its own arm} Rather than being excluded, inadmissibility is run
deliberately, in both refusal modes. This is what keeps the refusal channel from contaminating arms
that claim to measure placement: the channel is given somewhere to appear, and \cref{prop:capacity}
predicts exactly where. The padded refusal mode is the same arm with the fix applied, so the two
together measure the countermeasure rather than merely asserting it.

\end{document}
